% from main.tex
%%%%%%%%%%  scan idx 127  |  folio 114  %%%%%%%%%%
où les
flèches horizontales sont des épimorphismes et \(G'\in\Ob\mathcal C_0\). Comme \(\Ker(f)\) est quotient de \(\Ker(fp)\), il suffit de démontrer que \(\Ker(fp)\) est de type fini, donc on peut supposer que \(H'=G'\). Les flèches \(F\to G\) et \(H'\to G\) définissent \(g:F\oplus H'\to G\). Posant \(K=\Ker(G\to H)\), on a donc un diagramme commutatif, où les lignes sont exactes et \(q\) est un épimorphisme :
\[
\begin{tikzcd}
0\arrow[r]&F\arrow[r]\arrow[d,"q"']&F\oplus H'\arrow[r]\arrow[d,"g"]&H'\arrow[r]\arrow[d,"f"]&0\\
0\arrow[r]&K\arrow[r]&G\arrow[r]&H\arrow[r]&0 .
\end{tikzcd}
\]
Le diagramme du serpent montre que \(\Ker(g)\to\Ker(f)\) est un épimorphisme. Mais, comme \(F\oplus H'\) est de type fini et \(G\) précohérent, \(\Ker(g)\) est de type fini d'après a), et l'on gagne.

c) Prouvons d'abord l'assertion non respée. À un changement de notation près, il s'agit de montrer que si \(g:G'\to G\) est un morphisme de \(\mathcal C_S\) avec \(G'\in\Ob\mathcal C_0\), \(\Ker(g)\) est de type fini. Notant \(H'\) l'image de la flèche composée \(G'\to H\), on a un diagramme commutatif dont les lignes sont exactes
\[
\begin{tikzcd}
0\arrow[r]&F'\arrow[r]\arrow[d,"f"']&G'\arrow[r]\arrow[d,"g"]&H'\arrow[r]\arrow[d,"h"]&0\\
0\arrow[r]&F\arrow[r]&G\arrow[r]&H\arrow[r]&0 .
\end{tikzcd}
\]
Comme \(h\) est un monomorphisme, le diagramme du serpent donne un isomorphisme \(\Ker(f)\qis\Ker(g)\). Mais \(F'\) est de type fini d'après a), donc (à nouveau grâce à a)) \(\Ker(f)\) est de type fini, ce qui prouve l'assertion non respée. Passons à l'assertion respée (en toute rigueur on ne peut invoquer (2.5 b)) car on ne suppose pas que \(\mathcal C_0\) est localement quasi-stable par noyau d'épimorphisme). Quitte à localiser, il existe des épimorphismes \(F'\xrightarrow{u}F\), \(H'\xrightarrow{v}H\),
